% GENERATED FILE. Edit canonical lessons, then run npm run generate:books.
% canonical-source-hash: fnv1a64:4f3a7c86e2d4216a
% canonical-lessons: ZH-C303-qiming, ZH-C303-xiaoming, ZH-C303-youeryuan, ZH-C303-kaixue, ZH-C303-xinsheng

\chapter{A Name and a School}
\label{ch:zh-a-name-and-a-school}
\hlchaptermodality{303}

\begin{chapteropening}
\textbf{By the end of this chapter:} \emph{I can talk about naming a baby, a childhood name, nursery school and the start of term.}

Complete the last lesson of chapter 303: I can talk about naming a baby, a childhood name, nursery school and the start of term.
\end{chapteropening}

\section[\texorpdfstring{\zh{起名} (qǐmíng)}{起名 (qǐmíng)}]{\zh{起名} (qǐmíng) — to name a baby; to give a name}

\label{lesson:ZH-C303-qiming}



\begin{quote}
Before the new one: say the Mandarin for to bring back to life, then the Mandarin for a pair of twins.
\end{quote}

\subsection*{You'll want to know: \zh{起名}}
\textbf{\zh{起名}} — \emph{qǐmíng} — \textquotedblleft{}to name a baby; to give a name\textquotedblright{}.

\begin{quote}
\emph{gěi háizi qǐmíng bù róngyì} — Choosing a name for a child isn't easy.
\end{quote}

\begin{cousinweb}
\textbf{\zh{起名}} (\emph{qǐmíng}) is built from \zh{起} 'to raise' and \zh{名} 'a name', both inherited from Middle Chinese, \emph{khiX} and \emph{mjieng} (Guangyun, in Baxter's transcription), which glosses the first 'to rise' and 'to start'.
\end{cousinweb}

\begin{grammarlens}[title={What you've built}]
One more word for telling what happened.
\end{grammarlens}

\subsection*{Your turn}
\begin{itemize}
\raggedright
  \item \emph{Say it:} \emph{qǐmíng}
  \item \emph{Say it:} \emph{qǐmíng}, once more
  \item \emph{Say it:} \emph{gěi háizi qǐmíng bù róngyì} — Choosing a name for a child isn't easy.
  \item \emph{Recall:} say the Mandarin for to bring back to life, then the Mandarin for a pair of twins, then say \emph{qǐmíng} again
\end{itemize}

\begin{tcolorbox}[breakable,colback=teal!4,colframe=teal!35!black,title={Before you move on}]
What does \zh{起名} mean? (\textquotedblleft{}To name a baby; to give a name\textquotedblright{}.) Say it once more.
\end{tcolorbox}
\section[\texorpdfstring{\zh{小名} (xiǎomíng)}{小名 (xiǎomíng)}]{\zh{小名} (xiǎomíng) — childhood pet name}

\label{lesson:ZH-C303-xiaoming}



\begin{quote}
Before the new one: say the Mandarin for a pair of twins, then the Mandarin for to name a baby.
\end{quote}

\subsection*{You'll want to know: \zh{小名}}
\textbf{\zh{小名}} — \emph{xiǎomíng} — \textquotedblleft{}childhood pet name\textquotedblright{}.

\begin{quote}
\textbf{\zh{你的小名是什么}} \emph{nǐ de xiǎomíng shì shénme} — What was your pet name as a child?
\end{quote}

\begin{cousinweb}
\textbf{\zh{小名}} (\emph{xiǎomíng}) is built from \zh{小} 'small' and \zh{名} 'a name', both inherited from Middle Chinese, \emph{sjewX} and \emph{mjieng} (Guangyun, in Baxter's transcription): the small name, used at home.
\end{cousinweb}

\begin{grammarlens}[title={What you've built}]
One more word for telling what happened.
\end{grammarlens}

\subsection*{Your turn}
\begin{itemize}
\raggedright
  \item \emph{Say it:} \emph{xiǎomíng}
  \item \emph{Say it:} \emph{xiǎomíng}, once more
  \item \emph{Say it:} \emph{nǐ de xiǎomíng shì shénme} — What was your pet name as a child?
  \item \emph{Recall:} say the Mandarin for a pair of twins, then the Mandarin for to name a baby, then say \emph{xiǎomíng} again
\end{itemize}

\begin{tcolorbox}[breakable,colback=teal!4,colframe=teal!35!black,title={Before you move on}]
What does \zh{小名} mean? (\textquotedblleft{}Childhood pet name\textquotedblright{}.) Say it once more.
\end{tcolorbox}
\section[yòu'éryuán]{yòu'éryuán — kindergarten; nursery school}

\label{lesson:ZH-C303-youeryuan}



\begin{quote}
Before the new one: say the Mandarin for to name a baby, then the Mandarin for a childhood pet name.
\end{quote}

\subsection*{You'll want to know: yòu'éryuán}
\emph{yòu'éryuán} — \textquotedblleft{}kindergarten; nursery school\textquotedblright{}.

\begin{quote}
\emph{háizi sān suì de shíhou kāishǐ qù yòu'éryuán} — The child started kindergarten at three.
\end{quote}

Its characters come later; for now it is a word for the ear.

\begin{cousinweb}
\emph{yòu'éryuán} is built from \emph{yòu} 'young', \zh{儿} 'a child' and \emph{yuán} 'a garden', all inherited from Middle Chinese, \emph{'jiwH}, \emph{nye} and \emph{hjwon} (Guangyun, in Baxter's transcription), which glosses the first 'young' and the third 'a garden': a garden of young children.
\end{cousinweb}

\begin{grammarlens}[title={What you've built}]
One more word for telling what happened.
\end{grammarlens}

\subsection*{Your turn}
\begin{itemize}
\raggedright
  \item \emph{Say it:} \emph{yòu'éryuán}
  \item \emph{Say it:} \emph{yòu'éryuán}, once more
  \item \emph{Say it:} \emph{háizi sān suì de shíhou kāishǐ qù yòu'éryuán} — The child started kindergarten at three.
  \item \emph{Recall:} say the Mandarin for to name a baby, then the Mandarin for a childhood pet name, then say \emph{yòu'éryuán} again
\end{itemize}

\begin{tcolorbox}[breakable,colback=teal!4,colframe=teal!35!black,title={Before you move on}]
What does yòu'éryuán mean? (\textquotedblleft{}Kindergarten; nursery school\textquotedblright{}.) Say it once more.
\end{tcolorbox}
\section[\texorpdfstring{\zh{开学} (kāixué)}{开学 (kāixué)}]{\zh{开学} (kāixué) — to start the school term}

\label{lesson:ZH-C303-kaixue}



\begin{quote}
Before the new one: say the Mandarin for a childhood pet name, then the Mandarin for a kindergarten.
\end{quote}

\subsection*{You'll want to know: \zh{开学}}
\textbf{\zh{开学}} — \emph{kāixué} — \textquotedblleft{}to start the school term\textquotedblright{}.

\begin{quote}
\emph{kāixué de shíhou, wǒ rènshi le tā} — I got to know him when term started.
\end{quote}

\begin{cousinweb}
\textbf{\zh{开学}} (\emph{kāixué}) is built from \zh{开} 'to open' and \zh{学} 'to study', both inherited from Middle Chinese, \emph{khoj} and \emph{haewk} (Guangyun, in Baxter's transcription), which glosses the first 'to open up': the opening of the school year.
\end{cousinweb}

\begin{grammarlens}[title={What you've built}]
One more word for telling what happened.
\end{grammarlens}

\subsection*{Your turn}
\begin{itemize}
\raggedright
  \item \emph{Say it:} \emph{kāixué}
  \item \emph{Say it:} \emph{kāixué}, once more
  \item \emph{Say it:} \emph{kāixué de shíhou, wǒ rènshi le tā} — I got to know him when term started.
  \item \emph{Recall:} say the Mandarin for a childhood pet name, then the Mandarin for a kindergarten, then say \emph{kāixué} again
\end{itemize}

\begin{tcolorbox}[breakable,colback=teal!4,colframe=teal!35!black,title={Before you move on}]
What does \zh{开学} mean? (\textquotedblleft{}To start the school term\textquotedblright{}.) Say it once more.
\end{tcolorbox}
\section[\texorpdfstring{\zh{新生} (xīnshēng)}{新生 (xīnshēng)}]{\zh{新生} (xīnshēng) — new student}

\label{lesson:ZH-C303-xinsheng}



\begin{quote}
Before the new one: say the Mandarin for a kindergarten, then the Mandarin for to start the school term.
\end{quote}

\subsection*{You'll want to know: \zh{新生}}
\textbf{\zh{新生}} — \emph{xīnshēng} — \textquotedblleft{}new student\textquotedblright{}.

\begin{quote}
\textbf{\zh{新生下星期开学}} \emph{xīnshēng xià xīngqī kāixué} — The new students start next week.
\end{quote}

\begin{cousinweb}
\textbf{\zh{新生}} (\emph{xīnshēng}) is built from \zh{新} 'new' and \zh{生}, both inherited from Middle Chinese, \emph{sin} and \emph{srjaeng} (Guangyun, in Baxter's transcription), which glosses the second 'to grow'; its sense 'a student' is not given there.
\end{cousinweb}

\begin{grammarlens}[title={What you've built}]
One more word for telling what happened.
\end{grammarlens}

\subsection*{Your turn}
\begin{itemize}
\raggedright
  \item \emph{Say it:} \emph{xīnshēng}
  \item \emph{Say it:} \emph{xīnshēng}, once more
  \item \emph{Say it:} \emph{xīnshēng xià xīngqī kāixué} — The new students start next week.
  \item \emph{Recall:} say the Mandarin for a kindergarten, then the Mandarin for to start the school term, then say \emph{xīnshēng} again
\end{itemize}

\begin{tcolorbox}[breakable,colback=teal!4,colframe=teal!35!black,title={Before you move on}]
What does \zh{新生} mean? (\textquotedblleft{}New student\textquotedblright{}.) Say it once more.
\end{tcolorbox}
